\section{Experimental architecture}

\subsection{Persistent discovery archive}

Genesis stores evaluated organism variants in a persistent content-addressed archive. A node records its parent, lineage depth, public quality, whether it was ever champion, a research signature describing the mutation family and mechanisms, and the source archive needed for exact replay. Non-champion but viable nodes remain available as stepping stones.

The archive design draws on open-ended evolutionary systems such as the Darwin G\"odel Machine, which preserve multiple branches rather than greedily replacing a single incumbent \cite{zhang2025dgm}. Genesis adds a stricter separation between archive viability and champion promotion. A child can remain useful for future exploration without displacing the current champion.

\subsection{Replay worlds}

A replay world exposes only the portion of a previously recorded discovery tree that a policy could have known at each step. Selecting a historical parent either reveals its recorded child or consumes a represented request when no continuation was recorded. The policy never receives an oracle for unexecuted mutations. Replay therefore simulates allocation decisions over observed history rather than predicting counterfactual organism outcomes.

This matches the central Dream-RSI insight: historical discovery trees can provide low-cost feedback for alternative exploration policies because realized outcomes are already known \cite{zheng2026dreamrsi}. Genesis uses that idea conservatively. Replay may select a search policy for later deployment, but real organism retention is still decided by the external evaluator.

\subsection{External organism evaluator}

The hidden evaluator is fixed for the real-organism epoch. It executes candidate software in a sandbox, checks hard safeguards and behavior contracts, computes capability scores and quality, and decides whether a candidate is archive-viable or champion-eligible. Its implementation and hidden cases are outside the editable organism and policy surfaces.

\subsection{Prospective protocol discipline}

The experiments use four recurring controls:
\begin{enumerate}
  \item \textbf{freeze before observation}: success rules, policies, hashes and seed construction are committed before the relevant hidden or holdout result exists;
  \item \textbf{one-shot evaluation}: a completed scientific result cannot be rerolled under the same protocol;
  \item \textbf{negative evidence retention}: failed attempts remain in the evidence register and may motivate a new protocol but cannot be reclassified;
  \item \textbf{content addressing}: policies, proposal patches, manifests and result records are bound by SHA-256 digests.
\end{enumerate}
These controls do not make the project independent of its author, but they reduce the space for post-hoc selection within the recorded experimental line.
